\printconcepts

\exercise{T/F: Let $f$ be a position function. The average rate of change on $[a,b]$ is the slope of the line through the points $(a, f(a))$ and $(b,f(b))$.}{T}

% cut for parity
%\exercise{T/F: The definition of the derivative of a function at a point involves taking a limit.}{T}

\exercise{In your own words, explain the difference between the average rate of change and instantaneous rate of change.}{Answers will vary.}

\exercise{In your own words, explain the difference between Definitions \ref{def:derivative_at_a_point} and \ref{def:the_derivative}.}{Answers will vary.}

\exercise{Let $y=f(x)$. Give three different notations equivalent to ``$\fp(x)$.''}{Answers will vary.}

% in Hartman, but we omit normal lines
%\exercise{If two lines are perpendicular, what is true of their slopes?}{Their product is $-1$.}

\printproblems

\input{exercises/02-01-exset-02}

\input{exercises/02-01-exset-03}

\input{exercises/02-01-exset-01}

\exercise{The graph of $f(x)=x^2-1$ is shown. 
	\begin{enumext}
	\item		Use the graph to approximate the slope of the tangent line to $f$ at the following points: $(-1,0)$, $(0,-1)$ and $(2,3)$. 
	\item		Using the definition, find $\fp(x)$.
	\item		Find the slope of the tangent line at the points $(-1,0)$, $(0,-1)$ and $(2,3)$.
	\end{enumext}
\begin{tikzpicture}[alt={A U shaped graph through (-1,0), (0,-1), and (1,0).}]
\begin{axis}[width=\marginparwidth,tick label style={font=\scriptsize},
axis y line=middle,axis x line=middle,name=myplot,
ytick={-1,1,2,3},ymin=-1.3,ymax=3.5,xmin=-2.1,xmax=2.1,grid=major]
\addplot [thick,draw={\colorone},smooth,domain=-2.1:2.1] {x^2-1};
\end{axis}
\node [right] at (myplot.right of origin) {\scriptsize $x$};
\node [above] at (myplot.above origin) {\scriptsize $y$};
\end{tikzpicture}}{\mbox{}\\[-2\baselineskip]\parbox[t]{\linewidth}{\begin{enumext}
\item	Approximations will vary; they should match (c) closely.
\item	$\fp(x) = 2x$
\item	At $(-1,0)$, slope is $-2$. At $(0,-1)$, slope is 0. At $(2,3)$, slope is 4.
\end{enumext}}}

\exercise{The graph of $\ds f(x)=\frac{1}{x+1}$ is shown. 
	\begin{enumext}
	\item		Use the graph to approximate the slope of the tangent line to $f$ at the following points: $(0,1)$ and $(1,0.5)$. 
	\item		Using the definition, find $\fp(x)$.
	\item		Find the slope of the tangent line at the points $(0,1)$ and $(1,0.5)$.
	\end{enumext}
\begin{tikzpicture}[alt={A curve starting up near (-1,5) and then moving downward toward (3,0).}]
\begin{axis}[width=\marginparwidth,tick label style={font=\scriptsize},
axis y line=middle,axis x line=middle,name=myplot,ytick={1,2,3,4,5},
ymin=-.5,ymax=5.5,xmin=-1.1,xmax=3.1,grid=major]
\addplot [thick,draw={\colorone},smooth,domain=-.9:3.1,samples=50] {1/(x+1)};
\end{axis}
\node [right] at (myplot.right of origin) {\scriptsize $x$};
\node [above] at (myplot.above origin) {\scriptsize $y$};
\end{tikzpicture}}{\mbox{}\\[-2\baselineskip]\parbox[t]{\linewidth}{\begin{enumext}
\item	Approximations will vary; they should match (c) closely.
\item		$\fp(x) = -1/(x+1)^2$
\item		At $(0,1)$, slope is $-1$. At $(1,0.5)$, slope is $-1/4$.
\end{enumext}}}

\input{exercises/02-01-exset-04}

\exercisesetinstructions{, a graph of $g(x)$ is given.  Using the graph, answer the following questions.
\begin{enumerate}
\begin{multicols}{2}
	\item	Where is $g(x) > 0$?
	\item	Where is $g(x) < 0$?
	\item	Where is $g(x) = 0$?
	\item	Where is $g'(x) > 0$?
	\item	Where is $g'(x) < 0$?
	\item	Where is $g'(x) = 0$?
\end{multicols}
\end{enumerate}}

\exercise{\mbox{}\\[-2\baselineskip]\begin{tikzpicture}[alt={A curve starting near (-3,-5), passing through (-2,0), going up to (-1,1), turning downward to the origin and down to (1,-1), turning back upward to (2,0) and finishing near (3,5).}]
\begin{axis}[width=\marginparwidth,tick label style={font=\scriptsize},
axis y line=middle,axis x line=middle,name=myplot,xtick={-2,-1,1,2},
ymin=-5.5,ymax=5.5,xmin=-3.1,xmax=3.1]
\addplot [thick,draw={\colorone},domain=-3.1:3.1,samples=40] {(x-2)*x*(x+2)};
\end{axis}
\node [right] at (myplot.right of origin) {\scriptsize $x$};
\node [above] at (myplot.above origin) {\scriptsize $y$};
\end{tikzpicture}}{\mbox{}\\[-2\baselineskip]\parbox[t]{\linewidth}{\begin{enumext}
\item		Approximately on $(-2,0)$ and $(2,\infty)$.
\item		Approximately on $(-\infty,-2)$ and $(0,2)$.
\item		Approximately at $x=0,\ \pm 2$.
\item		Approximately on $(-\infty,-1)$ and $(1,\infty)$.
\item		Approximately on $(-1,1)$.
\item		Approximately at $x=\pm 1$.
\end{enumext}}}

\exercise{\mbox{}\\[-2\baselineskip]\begin{tikzpicture}[alt={A curve starting near (-2,-6), going through (-1.5,0), up to (-1,4), turning at that point downward to (0,1), turning at that point up to (1,4), turning back down to (1.5,0), and finishing near (2,-6).},scale=.8]
\begin{axis}[width=\marginparwidth,tick label style={font=\scriptsize},
axis y line=middle,axis x line=middle,name=myplot,xtick={-2,-1,1,2},
ymin=-7.9,ymax=5.9,xmin=-2.5,xmax=2.5]
\addplot [thick,draw={\colorone},smooth,domain=-2.3:2.3,samples=50] {(-10)*(x^4/4-x^2/2)+1};
\end{axis}
\node [right] at (myplot.right of origin) {\scriptsize $x$};
\node [above] at (myplot.above origin) {\scriptsize $y$};
\end{tikzpicture}}{\mbox{}\\[-2\baselineskip]\parbox[t]{\linewidth}{\begin{enumext}
	\item	Approximately on $(-1.5,1.5)$.
	\item	Approximately on $(-\infty,-1.5) \cup (1.5,\infty)$.
	\item	Approximately at $x=\pm 1.5$.
	\item	On $(-\infty,-1) \cup (0,1)$.
	\item	On $(-1,0) \cup (1,\infty)$.
	\item	At $x=\pm 1$ and $x=0$.
\end{enumext}}}

\exercisesetend

\exercise{Suppose\label{diffimpliescont} that $f(x)$ is defined on an open interval containing the number $c$. Suppose that the limit
\[\lim_{h\to0}\frac{f(c+h)-f(c)}h\]
exists. Show that $f(x)$ is continuous at $x=c$. This shows that we could drop the assumption that $f(x)$ is continuous at $x=c$ in the definition of $\fp(c)$.}{}

% these problems should be much later
%\exerciseset{In Exercises}{, a function $f(x)$ is given, along with its domain and derivative. Determine if $f(x)$ is differentiable on its domain.}{
%
%\exercise{$\ds f(x) = \sqrt{x^5(1-x)}$, domain = $[0,1]$, $\ds\fp(x) = \frac{(5-6x)x^{3/2}}{2\sqrt{1-x}}$}{$\lim_{h\to0^+}\frac{f(0+h)-f(0)}h=0$; note also that $\lim_{x\to0^+}\fp(x)=0$. So $f$ is differentiable at $x=0$.\\
%$\lim_{h\to0^-}\frac{f(1+h)-f(1)}h=-\infty$; note also that $\lim_{x\to1^-}\fp(x) = -\infty$. So $f$ is not differentiable at $x=1$.\\
%$f$ is differentiable on $[0,1)$, not its entire domain.}
%
%\exercise{$\ds f(x) = \cos\big(\sqrt{x}\big)$,\ domain = $[0,\infty)$,\ $\ds\fp(x) = -\frac{\sin\big(\sqrt{x}\big)}{2\sqrt{x}}$}{The limit of the difference quotient is difficult to evaluate. Using \autoref{thm:special_limits}, we can determine $\lim_{x\to0^+}\fp(x) = -1/2$.\\
%Since $\fp$ is defined on $(0,\infty)$, we conclude $f$ is differentiable on $[0,\infty)$.}
%
%}

\printreview

\exercise{Approximate $\ds \lim_{x\to 5}\frac{x^2+2 x-35}{x^2-10.5 x+27.5}$.}{Approximately $-24$.}

\exercise{Use the Bisection Method to approximate, accurate to two decimal places, the root of $g(x) = x^3+x^2+x-1$ on $[0.5,0.6]$.}{Approximately $0.54$.}

\exercise{Give intervals on which each of the following functions are continuous.
\begin{multicols}{2}
\begin{enumext}
\item		$\ds \frac{1}{e^x+1}$
\item		$\ds \frac{1}{x^2-1}$
\item		$\ds \sqrt{5-x}\phantom{\frac{1}{e^x+1}}$
\item		$\ds \sqrt{5-x^2}\phantom{\frac{1}{x^2-1}}$
\end{enumext}
\end{multicols}}{\mbox{}\\[-2\baselineskip]\parbox[t]{\linewidth}{\begin{enumext}
\item		$(-\infty,\infty)$
\item		$(-\infty,-1)\cup (-1,1) \cup (1,\infty)$
\item		$(-\infty,5]$
\item		$[-\sqrt5,\sqrt5]$
\end{enumext}}}

\exercise{Use the graph of $f(x)$ provided to answer the following.
\begin{multicols}{2}
\begin{enumext}
\item		$\ds \lim_{x\to-3^-} f(x) = $?
\item		$\ds \lim_{x\to-3^+} f(x) = $?
\item		$\ds \lim_{x\to-3} f(x) = $?
\item		Where is $f$ continuous?
\end{enumext}
\end{multicols}
\begin{tikzpicture}[alt={A curve starting near (-5,3) and finishing at a hollow dot at (-3,1), then a solid dot at (-3,2), and then another curve starting from (-3,3) and finishing near (-1,-1).}]
\begin{axis}[width=\marginparwidth,tick label style={font=\scriptsize},
axis y line=middle,axis x line=middle,name=myplot,
ymin=-1.1,ymax=3.1,xmin=-5.1,xmax=0.5]
\addplot [thick,draw={\colorone},smooth,domain=-5:-3] {(x+3)^2+1};
\addplot [thick,draw={\colorone},smooth,domain=-3:0.5] {-(x+3)^2+3};
\filldraw [fill=white] (axis cs:-3,1) circle (1.5pt);
\filldraw [fill=white] (axis cs:-3,3) circle (1.5pt);
\filldraw [fill=black] (axis cs:-3,2) circle (1.5pt);
\end{axis}
\node [right] at (myplot.right of origin) {\scriptsize $x$};
\node [above] at (myplot.above origin) {\scriptsize $y$};
\end{tikzpicture}}{\mbox{}\\[-2\baselineskip]\parbox[t]{\linewidth}{\begin{enumext}
\item		1
\item		3
\item		Does not exist
\item		$(-\infty,-3)\cup (3,\infty)$
\end{enumext}}}
